\section{Additive 2D Curve Fitting}
\label{sec:fitter}

This appendix describes the downstream postprocess.
It is separated from the body because it is not part of the model: it runs after
thresholding, changes no weights, and is measured on its own terms.
Keeping the boundary explicit is deliberate.
A gap-repair stage that is reported jointly with a segmenter can quietly become
a second segmentation method, and an improvement in one stage can then hide a
regression in the other.

\subsection{Method}

The postprocess repairs short omissions in the thresholded predictions.
It is additive only: accepted joins set foreground pixels but never erase model
output, which makes every intervention localizable and reversible.

For each $z$-slice, we apply Zhang--Suen thinning~\cite{zhang1984fast}, remove
short spurs, label the 8-connected skeleton and mask components, and extract
curve endpoints with outward tangents.
We then snap endpoints to the mask edge so distances describe the visible gap
rather than skeleton retraction.
Endpoints clipped by the volume or missing data are retained for audit but may
not be matched.

We score endpoint pairs using distance, facing, adjacent-plane evidence, and
cross-slice track support.
Before a pair enters deterministic greedy matching it must satisfy the
following requirements:

\begin{enumerate}
  \item it does not close the same skeleton or mask component;
  \item its tangents face the gap and are sufficiently opposed;
  \item its reach is at most 21 pixels, while the calibrated safe regime at six
  pixels or less requires less supporting evidence;
  \item if an umbilicus is available, its radial displacement is at most three
  pixels;
  \item its cubic Bezier has a bounded arc-to-chord ratio and does not approach
  a third component;
  \item the path corridor does not cross another sheet on nearby slices; and
  \item it neither crosses nor touches an already accepted join, and neither
  endpoint has taken a better match.
\end{enumerate}

A first matching pass seeds endpoint and connection tracks.
A second pass follows each new track through $z$, uses connection confidence,
and retains only joins supported across at least three slices.
We then paint a cubic Bezier disk stroke and record every added pixel.

\subsection{Measured Behaviour}

In a calibrated $4\!\times\!5\!\times\!5$ study the fitter accepted 481 joins
and added 12,697 pixels.
Registered high-resolution precision was 92.5\% overall and 98.4\% for joins no
longer than six pixels.
We observed no full-turn fusion and measured a 21\% reduction in atlas overlap.
A larger 21-cube census measured 98.3\% weighted precision.

Adjacent high-resolution connectivity is useful confirmation but is not, by
itself, proof that a join remains on one wrap.
The umbilicus-aware radial gate is our primary cross-wrap safeguard.
A release evaluation should therefore report false joins by gate, distance, and
radial regime rather than only an aggregate precision.

A bridge scanner reports thin-neck weld candidates, but cutting is disabled by
default because erasure breaks the additive contract.
An experimental mesh-stage slab splitter separated zero of 246 fused runs; it is
retained in the source tree as a measured negative result and is not part of the
recommended pipeline.

\subsection{Later Postprocessors}

Three further opt-in stages are implemented in the repository and evaluated in
their own records rather than here: a checked repair profile that recovers more
local continuations under a tighter radial guard, a probability-guided
continuity tool that reaches beyond the native fitter's 21-pixel limit, and a
stage that extends supported curves into explicit 3D patches.
All three share this appendix's contract---additive, opt-in, separately
measured, no model change---and all three currently pin the earlier operating
threshold, so each must be re-qualified before being combined with the shipped
inference recipe.
None is a production default, and none should be read as a global threshold
reduction or an anatomical certification.

\begin{figure*}[t]
  \centering
  \figureasset{figure_3_line_fitter_examples}{0.96\linewidth}{2.9in}
  \caption{Measured line-fitter behavior on the released M7 baseline. Left:
  representative M7 masks before and after additive accepted joins. Center:
  the same connection persists across neighboring $z$-slices. Right:
  registered fine-CT evidence distinguishes a supported connection from a
  required separation. The calibration accepted 481 of 772 proposals, added
  12,697 pixels, reached 92.5\% registered fine-CT precision, and produced no
  observed full-turn fusion. Cyan is released M7 and green is fitted addition.}
  \Description{Released M7 before-and-after curve fitting, cross-slice track
  persistence, connected and separate fine-CT validation examples, and measured
  calibration statistics.}
  \label{fig:fitter}
\end{figure*}
